\documentclass[
	classe=$2^{de}$,
	headerTitle=Activité,
	twocolumn,
	landscape
]{exercice}

\usetikzlibrary{positioning,calc}

\renewcommand{\arraystretch}{1.5}
\setlength{\columnsep}{1cm}

\title{Activité : évolution réciproque}

\begin{document}

\newcommand{\Activite}{
\maketitle

Une entreprise compte $250$ salariés en 2020.

En 2021, après une hausse considérable de son activité, elle décide d'embaucher $500$ personnes supplémentaires, pendant une durée d'un an.

\begin{enumerate}
	\item Quelle est la proportion des nouveaux employés dans l'entreprise en 2021 ? \correctionDots{$2/3$}
	\item Calculer le \textbf{coefficient} de cette évolution (c'est-à-dire, le nombre par lequel on a multiplié) : \correctionDots{$3$\hspace{1.5em}}.
	\item En 2022, le contrat des nouveaux arrivants arrive à son terme, et n'est pas renouvelé.

	      Quelle est l'évolution absolue du nombre d'employés ? \correctionDots{$-500$}

	      Quelle est l'évolution relative ? \correctionDots{$-2 / 3$}
	\item Remplir le schéma suivant :
	      \begin{center}
		      \begin{tikzpicture}[node distance=3cm]
			      \node[draw,circle] (D) {\correction{$250$}\hspace{2em}};
			      \node[text width=4cm] at ($(D) - (1.5,0)$) {\hspace{1em} employés\\ en 2020/2022};
			      \node[draw,circle] (A)[right=of D] {\correction{$750$}\hspace{2em}};
			      \node[text width=4cm] at ($(A) + (3,0)$) {employés\\ en 2021};
			      \path[->]
			      (D) edge[bend left] node[above]{$× \correctionDots{3}$} (A)
			      (A) edge[bend left] node[below]{$× \correctionDots{\dfrac{1}{3}}$} (D);
		      \end{tikzpicture}
	      \end{center}
\end{enumerate}
}

\Activite

\newpage
\Activite

\end{document}